\documentclass[10pt,a4paper]{amsart}
\usepackage[margin=19mm]{geometry}
\usepackage{amsmath,amssymb,mathtools}
\usepackage[scr=rsfs,cal=euler]{mathalfa}
\usepackage{xfrac}
\usepackage[unicode,hidelinks,bookmarks=false]{hyperref}
\newcommand{\ie}{i.e.\ }
\newcommand{\cf}{cf.\ }
\newcommand{\resp}{respectively\ }
\newcommand{\Subsection}{\subsection}
\providecommand{\nobreakdash}{\nobreak\hbox{-}}
\setlength{\emergencystretch}{2em}
\multlinegap0pt
\usepackage{amssymb}
\newtheorem{lemma}{Lemma}
\newcommand{\C}{\mathbb{C}}
\newcommand{\Gl}{{\mathrm{GL}}}
\newcommand\Mon{\mathrm{Mon}}
\newcommand{\PD}{\operatorname{PD}}
\newcommand{\R}{\mathbb{R}}
\newcommand{\Sl}{{\mathrm{SL}}}
\newcommand{\So}{{\mathrm{SO}}}
\newcommand{\Su}{{\mathrm{SU}}}
\newcommand{\U}{{\mathrm{U}}}
\newcommand{\Z}{\mathbb{Z}}
\newcommand{\cM}{\mathcal{M}}
\newcommand{\cT}{\mathcal{T}}
\title{Cancelling CR singularities of 3-manifolds in complex threefolds}
\author{Marko Slapar, Sa\v{s}o Strle}
\date{Source: arXiv:2606.21709v1 (CC BY 4.0)}
\begin{document}
\maketitle

\noindent\emph{Source and licence.} Cancelling CR singularities of 3-manifolds in complex threefolds. Version arXiv:2606.21709v1 (CC BY 4.0), distributed by arXiv under the Creative Commons Attribution 4.0 International licence, http://creativecommons.org/licenses/by/4.0/, as recorded in the arXiv licence metadata for this exact version and shown on the abstract page https://arxiv.org/abs/2606.21709v1. Full licence notice: \texttt{licenses/arxiv-2606.21709-CC-BY-4.0.txt}.\par\smallskip

\noindent\textit{Editorial scope.} Counted unit: the article's lemma lemma (source0.tex lines 222--236). The article's complete original proof of it is delivered with it (source0.tex lines 238--310, one \texttt{proof} environment of 73 lines). No mathematics is written, repaired or re-derived: the statement and the proof are the article's own lines, in the article's own notation, and no retained line is substituted or re-flowed.  No further source-local context is needed: every cross-reference of every retained line resolves inside the two delivered spans.  Nothing was omitted from any retained span, so the package carries no omission record.  No retained line contains a citation, so the wrapper carries no bibliography and no bibliography entry.  No retained line contains a figure environment, an image-inclusion command or drawing code, so no image is delivered and the package declares no asset dependency.  Editorial formatting only, in addition: an AMS \texttt{amsart} preamble that restates the article's own declarations and macros used by the retained lines, namely the environments \texttt{lemma} on separate counters, together with the standard packages those lines need.  The retained environments print in this document as Lemma 1 in order of appearance, whereas the article prints the same items with the numbering of its own full document.  The complete native source of the article as distributed in the arXiv e-print for this exact version is bundled unchanged as \texttt{sources/203287/source0.tex}.
\begin{lemma}\label{lemma}
Let $N$ be a compact oriented smooth $3$-manifold with boundary and let $E\to N$ be a complex vector bundle of rank $3$. Denote by 
\[\cM=\Mon_{\R}(TN,E)\] 
the bundle of $\R$-linear monomorphisms and let $\cT\subset \cM$ be the open subbundle consisting of totally real monomorphisms. Let
\[A\in \Gamma(N,\mathcal M)\]
be a section which takes values in $\cT$ over $\partial N$.  Choose a positive global frame $(e_1,e_2,e_3)$ of $TN$, and define
\[\Delta_A=A(e_1)\wedge_{\C} A(e_2)\wedge_{\C} A(e_3)
        \in \Gamma\left(N,\Lambda^3_{\C}E\right).\]
Assume that $\Delta_A$ is transverse to the zero section and let
\[Z_A=\Delta_A^{-1}(0),\]
with the orientation induced as the zero set of the section $\Delta_A$ of the complex line bundle $\Lambda^3_{\C}E$. 
Then $A$ is homotopic relative to $\partial N$, through sections of $\mathcal M$, to a section of $\mathcal T$ if and only if 
\[[Z_A]=0 \text{\ in\ } H_1(N;\Z).\]
Moreover, the homotopy may be chosen fixed on a sufficiently small collar of $\partial N$.
\end{lemma}

\begin{proof}
A real monomorphism \(B:T_pN\to E_p\) is totally real if and only if 
\[B(e_1)\wedge_{\C}B(e_2)\wedge_{\C}B(e_3)\ne 0.\] 
Thus the subbundle \(\cT\subset \cM\) is exactly the inverse image of the complement of the zero section under the determinant map. If a positive framing of $TN$ is changed by a map to $\Gl^+(3,\R)$, then $\Delta_A$ is multiplied by the positive real determinant of this change of frame. Hence the zero set $Z_A$, its induced orientation, and homology class $[Z_A]$ 
are independent of the chosen positive framing.

We apply obstruction theory to the bundle pair
\[(\cM,\cT)\longrightarrow N.\]
After choosing local positive frames of $TN$ and local complex frames of $E$, the fiber pair is identified with 
\[(V,W)=\left(\Mon_{\R}(\R^3,\C^3),\Gl(3,\C)\right).\]
The fiber $V$ deformation retracts to the real Stiefel manifold $V_{6,3}$, and $W$ deformation retracts to $\U(3)$. 
Since $V_{6,3}$ is $2$-connected, the relative homotopy exact sequence 
\begin{equation*}	
\begin{split}	
\cdots \longrightarrow
\pi_3(U(3))
\longrightarrow
\pi_3(V_{6,3})
\longrightarrow
\pi_3(V_{6,3},U(3))
\overset{\partial}{\longrightarrow}
\pi_2(U(3))
\longrightarrow
\pi_2(V_{6,3})\\
\longrightarrow
\pi_2(V_{6,3},U(3))
\overset{\partial}{\longrightarrow}
\pi_1(U(3))
\longrightarrow
\pi_1(V_{6,3})
\longrightarrow
\pi_1(V_{6,3},U(3))
\longrightarrow
0
\end{split}
\end{equation*}
gives
\[\pi_1(V,W)=0, \qquad \pi_2(V,W)\cong \pi_1(W)\cong \pi_1(\U(3))\cong \Z.\]
Under the connecting isomorphism $\pi_2(V,W)\cong \pi_1(W)$, the generator is detected by the determinant. Indeed,
\[\det\nolimits_{\C}:W=\Gl(3,\C)\longrightarrow \C^*\]
induces an isomorphism on $\pi_1$, since $\Sl (3,\C)\simeq \Su(3)$ is simply connected. Equivalently, the preferred generator of \(\pi_2(V,W)\) is represented by a small positively oriented normal disk to the codimension-two determinant locus $\{\det\nolimits_{\C}=0\}\subset V$; its boundary is a positive meridian.

There is no secondary obstruction in degree $3$. Indeed, the map 
\[\pi_3(W)\longrightarrow \pi_3(V)\]
is surjective. To see this, write $V\simeq \operatorname{SO}(6)/\So(3)$. The standard map $\U(3)\to \operatorname{SO}(6)$, obtained by forgetting the complex structure, is an isomorphism on $\pi_3$; for instance, this follows from the fibration
\[\U(3)\longrightarrow \operatorname{SO}(6)\longrightarrow \So(6)/U(3)
        \cong \C P^3\]
and the fact that $\pi_k(\mathbb{C}P^3)=0$ for $k=3,4$. The projection $\So(6)\to \So(6)/\So(3)$ is surjective on $\pi_3$, by the fibration
\[\So(3)\longrightarrow \So(6)\longrightarrow \So(6)/\So(3)\]
and $\pi_2(\operatorname{SO}(3))=0$. Hence
\[\pi_3(V,W)=0 .\]

We apply obstruction theory to address the question of when it is possible to deform $A$, relative to $\partial N$, through sections of $\cM$ to a section of $\cT$. Fixing a structure of a CW complex on $N$ it follows, using the triviality of $\pi_k(V,W)$ for $k=0, 1$, that $A$ is homotopic to a section that maps the 1-skeleton into $\cT$. Since for all $x\in N$, $\pi_2(\cM_x,\cT_x) \cong \pi_2(V,W) \cong \Z$, there is an obstruction class in 
\[H^2\left(N,\partial N;\underline{\pi_2(\cM_x,\cT_x)}\right)\]
whose vanishing is equivalent to existence of a section homotopic to $A$ that maps the 2-skeleton into $\cT$. If this obstruction vanishes, it follows from $\pi_3(V,W) = 0$ that $A$ can be homotoped, relative to $\partial N$, to a section of $\cT$.

Note that the local coefficient system $\underline{\pi_2(\cM_x,\cT_x)}$ is trivial. The determinant locus in $\cM_x$ is independent of the choice of trivializations of the bundles and hence defines a global codimension-two locus in $\cM$. Transporting the generator of $\pi_2(\cM_x,\cT_x)$, represented by a positively oriented disk normal to the determinant locus, around any loop will result in a positively oriented normal disk representing the same generator. This follows, since a change of positive frame in $TN$ and a change of complex frame in $E$ preserves the orientations. Alternatively, mapping the fibre pair by the determinant homomorphism $\det_{\C} : (V,W) \to (\C,\C^*)$ induces isomorphism on $\pi_2$. A change in frames multiplies the determinant by a nowhere-zero complex-valued function and hence preserves the generator of $\pi_2(\C,\C^*)$, showing that the local coefficient group $\underline{\pi_2\left(\Lambda^3_{\C}E_x,\Lambda^3_{\C}E_x\setminus\{0\}\right)}$ is trivial. It follows that the only obstruction to deforming $A$, relative to $\partial N$, through sections of $\cM$ into $\cT$ lies in
\[H^2(N,\partial N;\Z).\]

Next we identify this obstruction. By naturality of obstruction theory, the determinant map sends the primary obstruction for the pair $(\cM,\cT)$ to the primary obstruction for the pair
\[\left(\Lambda^3_{\C}E,\Lambda^3_{\C}E\setminus\{0\}\right).\]
Since the induced map
\[(\det\nolimits_\C)_*:\pi_2(V,W)\longrightarrow \pi_2(\mathbb C,\mathbb C^*)\]
is an isomorphism, these obstruction classes are identified. The latter obstruction is precisely the obstruction to deforming the determinant section $\Delta_A$, fixed on $\partial N$, to a nowhere-zero section of the complex line bundle $\Lambda^3_{\mathbb C}E$. Hence it is the relative first Chern class
\[c_1\left(\Lambda^3_{\mathbb C}E,\Delta_A|_{\partial N}\right)\in H^2(N,\partial N;\mathbb Z),\]
where the nonzero boundary section \(\Delta_A|_{\partial N}\) is used as the boundary trivialization. Since $\Delta_A$ is transverse to the zero section, the usual zero-locus representative of the relative Chern class gives
\[c_1\left(\Lambda^3_{\C}E,\Delta_A|_{\partial N}\right) =\PD_N[Z_A]\in H^2(N,\partial N;\Z),\]
with $Z_A$ oriented as the zero set of the section $\Delta_A$ of the complex line bundle $\Lambda^3_{\C}E$. Hence the obstruction vanishes if and only if
\[ [Z_A]=0.\]

Finally, since $\partial N$ is compact and $\cT$ is open, $A$ takes values in $\cT$ on a (smooth) collar of $\partial N$. 
Applying the same relative obstruction argument to the complement of this collar and keeping the new boundary fixed, and then extending the resulting homotopy by the constant homotopy on the collar, gives a homotopy fixed on a collar of $\partial N$.
\end{proof}

\end{document}
